\documentclass[11pt,a4paper]{article}
\usepackage{shared/parscover}

% ============================================================================
% Staking and Validator Economics for Pars Network
% Pars Network Technical Whitepaper PNP-013
% ============================================================================

\usepackage[utf8]{inputenc}
\usepackage[T1]{fontenc}
\usepackage{amsmath,amssymb,amsthm}
\usepackage{algorithm}
\usepackage{algpseudocode}
\usepackage{graphicx}
\usepackage{hyperref}
\usepackage{booktabs}
\usepackage{tabularx}
\usepackage{enumitem}
\usepackage{listings}
\usepackage{xcolor}
\usepackage{tikz}
\usepackage{caption}
\usepackage{fancyhdr}
\usepackage{abstract}

\geometry{margin=1in}

\hypersetup{
    colorlinks=true,
    linkcolor=blue!70!black,
    citecolor=green!50!black,
    urlcolor=blue!60!black
}

\lstset{
    basicstyle=\ttfamily\small,
    breaklines=true,
    frame=single,
    backgroundcolor=\color{gray!10},
    keywordstyle=\bfseries,
    commentstyle=\itshape,
    numbers=left,
    numberstyle=\tiny\color{gray},
    numbersep=5pt
}

\usetikzlibrary{arrows.meta,positioning,shapes.geometric,fit,calc}

\newtheorem{definition}{Definition}[section]
\newtheorem{theorem}{Theorem}[section]
\newtheorem{lemma}{Lemma}[section]
\newtheorem{proposition}{Proposition}[section]
\newtheorem{corollary}{Corollary}[section]

\pagestyle{fancy}
\fancyhf{}
\fancyhead[L]{\small Zach Kelling, Pars}
\fancyhead[R]{\small Staking Economics}
\fancyfoot[C]{\thepage}

% ============================================================================
\title{%
    \textbf{Staking and Validator Economics\\
    for Pars Network}\\[0.5em]
    \large Pars Network Technical Whitepaper PNP-013
}

\author{%
    Zach Kelling\\
    Pars DAO\\
    \texttt{research@pars.network}\\[0.5em]
    \small Version 1.0
}

\date{2026}

% ============================================================================
\begin{document}
\parscoverpage

\thispagestyle{fancy}

% ============================================================================
\begin{abstract}
\noindent
We present the staking and validator economics for the Pars Network,
extending the VEASHA tokenomics framework with concrete staking
mechanisms, validator incentives, slashing conditions, delegation, and
privacy-preserving stake proofs.  Pars validators secure the
privacy-preserving consensus layer (PNP-012) while maintaining
anonymity---a requirement that conflicts with traditional staking designs
where validator stakes are publicly visible.  We introduce
zero-knowledge stake proofs that demonstrate sufficient stake for
validator eligibility without revealing the exact amount or the
staker's identity.  The economic model incentivises geographic
distribution across diaspora jurisdictions, penalises correlation
(validators in the same jurisdiction or hosting provider), and provides
a delegation mechanism that preserves delegator privacy.  Slashing
conditions are designed for the coercion-resistant threat model: a
validator under duress can trigger panic mode without being slashed for
the resulting protocol deviation.
\end{abstract}

\vspace{1em}
\noindent\textbf{Keywords:} staking, validator economics, VEASHA,
slashing, delegation, zero-knowledge stake proofs, geographic
distribution, coercion-resistant slashing

\tableofcontents
\newpage

% ============================================================================
\section{Introduction}
\label{sec:intro}

Proof-of-stake networks secure consensus through economic bonds:
validators lock tokens as collateral, earning rewards for honest
participation and risking slashing for protocol violations.  The Pars
Network's privacy requirements introduce unique constraints on staking
design:

\begin{enumerate}
    \item \textbf{Stake privacy.}  Public stake amounts create power
    maps that adversaries use to identify high-value targets for
    coercion.  Validator stakes must be hidden.
    \item \textbf{Identity privacy.}  Validators participate
    anonymously (PNP-012).  Staking must not break anonymity.
    \item \textbf{Coercion tolerance.}  Slashing rules must
    accommodate panic-mode behaviour (PNP-005) without penalising
    coerced validators.
    \item \textbf{Geographic distribution.}  The diaspora is distributed
    across 60+ countries; staking incentives must encourage this
    distribution rather than concentrating in low-cost hosting regions.
\end{enumerate}

\subsection{Relation to VEASHA Tokenomics}

The VEASHA (Vote-Escrow ASHA) tokenomics paper defines the token
distribution, emission schedule, and governance voting power.  This
paper extends VEASHA with:

\begin{itemize}
    \item Concrete staking and unstaking mechanics.
    \item Validator reward distribution.
    \item Slashing conditions and dispute resolution.
    \item Delegation with privacy preservation.
    \item Zero-knowledge stake eligibility proofs.
\end{itemize}

\subsection{Contributions}

\begin{itemize}
    \item Zero-knowledge stake proofs for private validator eligibility.
    \item Correlation-penalised rewards encouraging geographic diversity.
    \item Coercion-aware slashing with panic-mode exceptions.
    \item Private delegation via stealth staking addresses.
    \item Formal analysis of economic security under the Pars threat model.
\end{itemize}

% ============================================================================
\section{VEASHA Staking}
\label{sec:veasha}

\subsection{Token Mechanics}

ASHA is the native token of the Pars Network.  VEASHA (vote-escrow ASHA)
is obtained by locking ASHA for a period $T$:

\begin{definition}[VEASHA Balance]
    A user locking $a$ ASHA for period $T$ (in weeks) receives VEASHA
    balance:
    \[
        \mathsf{veASHA}(a, T) = a \cdot \frac{T}{T_{\max}}
    \]
    where $T_{\max} = 208$ weeks (4 years).  VEASHA balance decays
    linearly as the lock approaches expiry.
\end{definition}

\subsection{Validator Registration}

To become a validator, a participant must:

\begin{enumerate}
    \item Lock at least $S_{\min}$ ASHA as VEASHA, where $S_{\min}$ is
    the minimum stake (initially 10{,}000 ASHA).
    \item Generate a validator key pair and register the public key on
    the Pars L2 staking contract.
    \item Generate a ZK proof of stake eligibility
    (Section~\ref{sec:zkstake}).
    \item Begin participating in consensus (PNP-012).
\end{enumerate}

The registration transaction links the VEASHA position to the validator
key via a commitment, without revealing the stake amount.

% ============================================================================
\section{Zero-Knowledge Stake Proofs}
\label{sec:zkstake}

\subsection{Motivation}

Traditional proof-of-stake systems publish validator stakes on-chain,
enabling stake-weighted leader election and reward distribution.  This
transparency is a liability for Pars: an adversary monitoring the
staking contract can identify the largest stakers and target them for
coercion.

\subsection{Construction}

We use a Groth16 proof (infrastructure from PNP-006) to prove stake
eligibility without revealing the amount.

\begin{definition}[Stake Eligibility Proof]
    Given a VEASHA commitment $C = \mathsf{Commit}(a, T, r)$ on-chain,
    the validator produces a ZK proof $\pi$ of the statement:
    \[
        \exists\, a, T, r : C = \mathsf{Commit}(a, T, r)
        \;\wedge\; a \cdot T / T_{\max} \geq S_{\min}
    \]
    The proof reveals nothing beyond the fact that the committed VEASHA
    balance exceeds $S_{\min}$.
\end{definition}

\subsection{Stake Ranges for Weighted Sampling}

Consensus (PNP-012) requires stake-weighted validator sampling.
Revealing exact weights breaks privacy.  Instead, we define stake
ranges:

\begin{table}[h]
\centering
\caption{Stake ranges for weighted consensus sampling.}
\label{tab:ranges}
\begin{tabular}{@{}lrr@{}}
\toprule
\textbf{Tier} & \textbf{VEASHA Range} & \textbf{Sampling Weight} \\
\midrule
Bronze  & $[S_{\min}, 3 S_{\min})$ & 1 \\
Silver  & $[3 S_{\min}, 10 S_{\min})$ & 3 \\
Gold    & $[10 S_{\min}, 30 S_{\min})$ & 8 \\
Diamond & $\geq 30 S_{\min}$ & 15 \\
\bottomrule
\end{tabular}
\end{table}

Validators prove membership in a tier via a range proof, revealing only
the tier, not the exact amount.

\begin{theorem}[Tier Privacy]
    Under the ZK proof soundness and zero-knowledge properties, the
    adversary learns only the validator's tier, not their exact stake.
    Within a tier, all validators are indistinguishable.
\end{theorem}

% ============================================================================
\section{Validator Rewards}
\label{sec:rewards}

\subsection{Emission Schedule}

Block rewards follow the VEASHA emission schedule: a decreasing emission
rate starting at 5\% annual inflation, halving every 4 years, with a
hard cap of 1 billion ASHA.

\begin{definition}[Block Reward]
    The reward for block $b$ at epoch $e$ is:
    \[
        R(b, e) = R_0 \cdot 2^{-e/E_{\mathrm{half}}} \cdot \frac{1}{B_e}
    \]
    where $R_0$ is the initial annual emission, $E_{\mathrm{half}}$ is
    the halving period in epochs, and $B_e$ is the number of blocks in
    epoch $e$.
\end{definition}

\subsection{Reward Distribution}

Rewards are distributed proportionally to sampling weight, adjusted by
two factors:

\paragraph{Uptime factor.}
\[
    f_{\mathrm{up}}(V_i) = \frac{\text{blocks participated}}{
    \text{blocks eligible}} \cdot \mathbf{1}[
    f_{\mathrm{up}} \geq 0.9]
\]
Validators below 90\% uptime receive zero rewards for that epoch.

\paragraph{Diversity factor.}
\[
    f_{\mathrm{div}}(V_i) = 1 - \alpha \cdot \frac{
    n_{\mathrm{same}}(V_i)}{n_{\mathrm{total}}}
\]
where $n_{\mathrm{same}}(V_i)$ is the number of validators sharing
the same jurisdiction or AS number as $V_i$, and $\alpha$ is a
correlation penalty parameter (initially $\alpha = 0.5$).

\begin{definition}[Validator Reward]
    Validator $V_i$ with tier weight $w_i$, uptime factor
    $f_{\mathrm{up}}(V_i)$, and diversity factor $f_{\mathrm{div}}(V_i)$
    receives epoch reward:
    \[
        r_i = R_{\mathrm{epoch}} \cdot
        \frac{w_i \cdot f_{\mathrm{up}}(V_i) \cdot f_{\mathrm{div}}(V_i)}{
        \sum_j w_j \cdot f_{\mathrm{up}}(V_j) \cdot f_{\mathrm{div}}(V_j)}
    \]
\end{definition}

\subsection{Geographic Distribution Incentive}

The diversity factor penalises validator concentration.  If 30\% of
validators are in a single jurisdiction, each receives a 15\% reward
reduction ($\alpha = 0.5$, $n_{\mathrm{same}}/n_{\mathrm{total}} = 0.3$).
This incentivises validators to distribute across diaspora hubs.

\begin{proposition}[Nash Equilibrium]
    Under the diversity-adjusted reward function, the Nash equilibrium
    validator distribution approximates the uniform distribution over
    available jurisdictions, weighted by the diaspora population
    distribution.
\end{proposition}

% ============================================================================
\section{Slashing}
\label{sec:slashing}

\subsection{Slashable Offences}

\begin{table}[h]
\centering
\caption{Slashing conditions and penalties.}
\label{tab:slashing}
\begin{tabular}{@{}llr@{}}
\toprule
\textbf{Offence} & \textbf{Evidence} & \textbf{Penalty} \\
\midrule
Double signing   & Two signatures at same height & 5\% \\
Prolonged downtime ($>$48h) & Missing from $>$48h of blocks & 1\% \\
Equivocation     & Conflicting votes in same round & 10\% \\
Censorship (proven) & Inclusion proof violation & 3\% \\
Correlated failure & $>$1/3 same-AS validators fail & 2$\times$ base \\
\bottomrule
\end{tabular}
\end{table}

\subsection{Coercion-Aware Slashing}

Standard slashing penalises protocol deviations indiscriminately.  Under
the Pars threat model, a coerced validator may deviate from the protocol
involuntarily.

\begin{definition}[Panic-Mode Exception]
    A validator in panic mode (PNP-005) whose protocol deviation matches
    the panic-mode behaviour profile is exempt from slashing.  The panic
    mode produces:
    \begin{enumerate}[label=(\roman*)]
        \item Valid signatures under the decoy key (not equivocation).
        \item Graceful withdrawal from consensus (not prolonged downtime
        if within 24h).
        \item Alert to the guardian network (PNP-010) for key resharing.
    \end{enumerate}
\end{definition}

\begin{theorem}[Coercion Safety]
    A coerced validator who activates panic mode is not slashed, provided:
    (a)~the panic credential is valid, (b)~the deviation matches the
    panic behaviour profile, and (c)~key resharing completes within 24h.
\end{theorem}

\subsection{Correlated Slashing}

Correlated failures---multiple validators failing simultaneously due to
shared infrastructure---indicate systemic risk.  Pars applies a
superlinear penalty:

\[
    \mathsf{slash}(V_i) = \mathsf{base\_penalty} \cdot
    \min\!\left(3,\; 1 + \frac{|\{V_j : \text{correlated with } V_i\}|}{
    n/3}\right)
\]

This incentivises infrastructure diversity: validators sharing a hosting
provider or jurisdiction face amplified slashing risk.

% ============================================================================
\section{Delegation}
\label{sec:delegation}

\subsection{Privacy-Preserving Delegation}

ASHA holders who do not wish to run validators can delegate their stake.
Delegation must preserve delegator privacy: the delegator's identity,
delegation amount, and choice of validator should not be publicly
visible.

\begin{definition}[Stealth Delegation]
    A delegator sends ASHA to a stealth staking address derived from
    the validator's public staking key, using the stealth address
    mechanism of PNP-007:
    \begin{enumerate}[label=(\roman*)]
        \item Delegator generates ephemeral key $r$ and computes
        stealth address $A_s = H(r \cdot V_{\mathrm{pub}}) \cdot G +
        V_{\mathrm{stake}}$.
        \item Delegator sends ASHA to $A_s$ with a VEASHA lock.
        \item Validator detects the delegation by scanning with their
        staking key.
        \item The delegation is invisible to external observers.
    \end{enumerate}
\end{definition}

\subsection{Reward Distribution to Delegators}

Validators distribute rewards to delegators via stealth payments:

\begin{enumerate}
    \item The validator computes each delegator's share based on their
    delegation weight.
    \item Rewards are sent to fresh stealth addresses derived from each
    delegator's viewing key.
    \item The on-chain observer sees only transfers between unlinkable
    addresses.
\end{enumerate}

\subsection{Delegation Commission}

Validators set a commission rate $c \in [0, 0.20]$ (max 20\%).  The
commission is enforced by the staking contract:

\[
    r_{\mathrm{delegator}} = r_{\mathrm{total}} \cdot
    \frac{d_i}{\sum_j d_j + s_{\mathrm{self}}} \cdot (1 - c)
\]

where $d_i$ is the delegator's stake and $s_{\mathrm{self}}$ is the
validator's self-stake.

% ============================================================================
\section{Unstaking and Withdrawal}
\label{sec:unstaking}

\subsection{Unbonding Period}

Unstaking requires an unbonding period of 14 days.  During unbonding:
\begin{itemize}
    \item The stake remains subject to slashing for offences committed
    before the unstaking request.
    \item The validator is removed from the active set.
    \item No rewards accrue.
\end{itemize}

\subsection{Emergency Unstaking}

Under the coercion model, a validator under threat may need to unstake
immediately.  Emergency unstaking is permitted with:
\begin{itemize}
    \item A 5\% early withdrawal penalty (burned, not redistributed).
    \item A valid panic credential (PNP-005).
    \item Guardian network confirmation (PNP-010).
\end{itemize}

% ============================================================================
\section{Economic Security Analysis}
\label{sec:security}

\subsection{Cost of Attack}

\begin{definition}[Byzantine Threshold Cost]
    The cost for an adversary to control $f = \lfloor (n-1)/3 \rfloor$
    validators is:
    \[
        C_{\mathrm{attack}} = f \cdot S_{\min} \cdot P_{\mathrm{ASHA}}
        + C_{\mathrm{infra}}(f)
    \]
    where $P_{\mathrm{ASHA}}$ is the market price of ASHA and
    $C_{\mathrm{infra}}$ is infrastructure cost.
\end{definition}

With $n = 100$ validators and $S_{\min} = 10{,}000$ ASHA, the
adversary must acquire and lock at least 330{,}000 ASHA, plus maintain
33 geographically distributed validator nodes.

\subsection{Stake Concentration Bound}

\begin{theorem}[Diversity Equilibrium]
    Under the diversity-adjusted reward function with $\alpha = 0.5$,
    no rational validator operator controls more than
    $\lceil n / J \rceil$ validators, where $J$ is the number of
    available jurisdictions.
\end{theorem}

\begin{proof}
    An operator adding a validator in a jurisdiction where they already
    operate $k$ validators receives reward factor
    $1 - 0.5 \cdot (k+1) / n_{\mathrm{local}}$.  When $k$ exceeds the
    jurisdiction's proportional share, the marginal reward becomes
    negative relative to operating in an unoccupied jurisdiction.
\end{proof}

\subsection{Long-Range Attack Resistance}

VEASHA's time-locked staking prevents long-range attacks: an adversary
cannot unstake, sell tokens, and then attempt to rewrite history with
worthless keys.  The 14-day unbonding period exceeds the finality
horizon.

% ============================================================================
\section{Governance Integration}
\label{sec:governance}

VEASHA balances determine governance voting power as defined in the
VEASHA tokenomics paper.  Staking extends this:

\begin{itemize}
    \item \textbf{Validator governance.}  Validators have additional
    governance rights: proposing protocol upgrades, adjusting staking
    parameters ($S_{\min}$, $\alpha$, commission caps).
    \item \textbf{Delegator governance.}  Delegators retain governance
    voting power proportional to their delegation, independent of the
    validator's vote.
    \item \textbf{Slashing governance.}  Disputed slashing events are
    resolved by VEASHA-weighted governance vote via the DAO (PNP-001
    governance).
\end{itemize}

% ============================================================================
\section{Related Work}
\label{sec:related}

\paragraph{Ethereum PoS.}  Ethereum's proof-of-stake uses 32 ETH minimum
stake with public validator identities.  Pars differs in requiring
validator anonymity and private stake amounts.

\paragraph{Penumbra.}  Penumbra~\cite{penumbra2023} provides private
staking on a Cosmos chain.  Pars adopts a similar ZK stake proof
approach but extends it with coercion-aware slashing and geographic
diversity incentives.

\paragraph{Lux Staking.}  The Lux network's staking provides the
validator registry infrastructure.  Pars validators are registered as
a subnet on Lux, inheriting the P-chain's validator management with
Pars-specific privacy extensions.

\paragraph{Curve veTokenomics.}  Curve's vote-escrowed CRV model
inspired VEASHA.  Pars extends the model with ZK proofs over the lock
position and staking-specific reward adjustments.

% ============================================================================
\section{Conclusion}
\label{sec:conclusion}

We have presented the staking and validator economics for the Pars
Network, extending the VEASHA tokenomics with concrete mechanisms for
private staking, validator rewards, coercion-aware slashing, and
privacy-preserving delegation.  Zero-knowledge stake proofs enable
validator eligibility verification without revealing stake amounts,
while geographic diversity incentives encourage validator distribution
across the diaspora.  The coercion-aware slashing framework
accommodates the unique threat model of communities operating under
authoritarian surveillance, ensuring that validators under duress are
not economically punished for activating panic protocols.

% ============================================================================
\begin{thebibliography}{99}

\bibitem{buterin2020pos}
V.~Buterin et al.,
``Combining GHOST and Casper,''
\emph{arXiv:2003.03052}, 2020.

\bibitem{penumbra2023}
Penumbra Labs,
``Penumbra: A Private Proof-of-Stake Network,''
2023.

\bibitem{curve2020ve}
Curve Finance,
``Curve DAO: Vote-Escrowed CRV,''
2020.

\bibitem{veasha}
Cyrus Pahlavi, Pars Team,
``VEASHA Tokenomics: Vote-Escrow Economic Model for Pars Network,''
2026.

\bibitem{pnp005}
Cyrus Pahlavi, Pars Team,
``PNP-005: Coercion-Resistant Protocols for Governance Under
Authoritarian Surveillance,''
2026.

\bibitem{pnp006}
Cyrus Pahlavi, Pars Team,
``PNP-006: Zero-Knowledge Identity: Privacy-Preserving Authentication,''
2026.

\bibitem{pnp007}
Cyrus Pahlavi, Pars Team,
``PNP-007: Stealth Addresses for Pars Network,''
2026.

\bibitem{pnp010}
Cyrus Pahlavi, Pars Team,
``PNP-010: Threshold Signatures for Pars Network Infrastructure,''
2026.

\bibitem{pnp012}
Cyrus Pahlavi, Pars Team,
``PNP-012: Privacy-Preserving Consensus for Pars Network,''
2026.

\bibitem{pars_dao}
Cyrus Pahlavi, Pars Team,
``Pars DAO Governance: Decentralised Autonomous Organisation for the
Persian Diaspora,''
2026.

\end{thebibliography}

\end{document}
